\chapter{电磁感应}
\section{电磁感应定律}
在奥斯特发现电流的磁效应之后, 人们致力于寻找其逆效应. 然而, 由于当时人们将探索的边界限制于静止的磁场, 于是迟迟未能实现突破. 直到法拉第用若干实验发现了电磁感应现象, 并且由后人总结出其数学形式, 形成了今天我们所见到的电磁感应定律.

在讨论电磁感应之前, 我们先来看一下电动势. 在一个电路中, 当我们的电流在整个电路中不相同时, 在某些地方会产生电荷的积累, 这些积累的电荷产生的电场将会阻碍电流的流动, 直到阻碍流动的电场与促进流动的电场到达平衡, 从而建立起整个电路的平衡. 也就是说, 在驱动电流在整个电路中流动的过程中涉及到电源的非静电力和静电力. 描述电池做功的能力时, 我们常常用非静电力在整个回路中的线积分来描述($\mathscr C$也可以写作$\varepsilon$)\[
\mathscr C = \oint_L  \mathbf{f}_s \cdot \d \mathbf{l}
\]
其中$\mathbf f_s$是单位电荷所受的非静电力

在理想电动势源中, 电荷受到的合力为$0$. 于是$E=-f_s$. 从而两板间的电势差为:
\[
V=-\int_a^b \mathbf{E} \cdot \d \mathbf{l} = \int_a^b f_s \cdot \d \mathbf{l} = \mathscr C
\]

法拉第曾经做了四个实验来验证他的猜想. 我们简单总结一下:
\begin{itemize}
  \item 取一个与电流计相连的线圈, 在其中插入和拔出磁棒时, 电流计中有电流流过.
  \item 将上面的磁棒换为载流线圈, 插入和拔出时, 电流计中有电流流过.
  \item 将载流线圈插入线圈, 控制开关使得载流线圈断电与通电, 开关断开和闭合后电流计中有电流流过.
  \item 取一个平面线框, 使其部分位于磁场区域内; 移动线框使穿过它的磁通量变化, 其中便有感应电流.
\end{itemize}

在这四个实验中, 我们发现其共性是通过导线回路或者线圈的磁通量发生了变化. 而且磁通量变化的越快, 电流越大.

精确的实验表明:

\begin{theorem}[法拉第电磁感应定律]
导体回路中感应电动势的大小与穿过回路磁通量的变化率成正比.

\[
\mathscr C =\oint_L \mathbf{f}_s\cdot\d\mathbf l =-\frac{\d \Phi_B}{\d t} =-\frac{\d}{\d t}\iint_S \mathbf{B} \cdot \d \mathbf{S}
\]

对于多匝线圈, 我们需要将每一匝线圈的磁通都计算, 若每匝线圈磁通都相同, 有:

\[
\mathscr C = -\frac{\d \Psi}{\d t} =-N\frac{\d \Phi_B}{\d t},\qquad \Psi=N\Phi_B
\]
\end{theorem}
其中$\Psi$被称为总磁通, 也称磁通匝链数或磁链.


关于如何判断感应电流的方向, 有楞次定律:
\begin{theorem}[楞次定律]
  闭合回路中感应电流的方向, 总是使得它所激发的磁场来阻碍引起感应电流的磁通量的变化.
\end{theorem}
楞次定律其实是能量守恒的必然结果. 感应电流在电路中流动时将释放焦耳热. 这部分能量必然来源于克服斥力或者引力所做的机械功. 否则插入磁棒的过程中感应电流既对外做功, 又释放焦耳热, 这显然是违反能量守恒的.

\section{动生电动势和感生电动势}
我们可以将产生电动势的原因简单分为两类. 一种是在恒磁场中因为导体运动而产生感应电动势. 另一种是导体不动, 因为磁场的变化而产生感应电动势. 前者被称为动生电动势, 后者被称为感生电动势.

先介绍动生电动势. 动生电动势可以被看作是洛伦兹力引起的.

\insertfig{4-1.png}{动生电动势与洛伦兹力}{0.25}

当导体以速度$v$向右运动时, 导体内的自由电子也以速度$v$跟随运动. 洛伦兹力的方向如图所示. 则在洛伦兹力的推动下, 自由电子将沿着$DCBA$方向运动, 从而电流沿着$ABCD$方向. 如果没有导体框与$CD$接触, 那么自由电子将向$C$端聚集, 从而使得$C$带负电, $D$带正电. 此时我们可以把$C$端看成负极, $D$看成正极. 这其中的非静电力就是作用在单位电荷上的洛伦兹力$\mathbf{v}\times\mathbf{B}$. 于是电动势为:
\[
\mathscr C = \int_C^D (\mathbf{v}\times \mathbf{B})\cdot \d \mathbf{l}
\]

上面这个公式对于任意的一段导体都是成立的.

事实上, 我们可以对任意的一个线圈$L$计算电动势, 上面的推导仍然适用. $L$可以闭合也可以不闭合, 于是整个线圈中的电动势为:
\[
\mathscr C = \int_L (\mathbf{v}\times \mathbf{B})\cdot \d \mathbf{l}
\]
我们给出下面的证明. 对于线圈上的一个线元$\d \mb{l}$, 其在$\d t$时间内扫过的面积为$\d \mb{l} \times \mathbf{v} \d t$.

考虑原有的线圈所围成的曲面$S$和$\d t$时间后线圈所围成的曲面$S'$. 这两个面和线圈扫过的侧面构成一整个闭合曲面. 考虑磁场的高斯定理:
\[
\oiint \mb{B} \cdot \d \mathbf{S} =
-\iint_S \mb{B} \cdot \d \mathbf{S} +\iint_S' \mb{B} \cdot \d \mathbf{S}
+\oint_{L} \mb{B} \cdot (\d \mb{l} \times \mathbf{v}) \d t = 0
\]

前两项相减得到$-\d \Phi$, 移项得到:
\[
-\d \Phi = -\oint_L \mathbf{B}\cdot (\d \mb{l}\times \mathbf{v}) \d t = \oint_L (\mathbf{v}\times \mathbf{B})\cdot \d \mathbf{l} \d t
\]
两边同时除以$\d t$

\[
\mathscr C = -\frac{\d \Phi}{\d t} =  \oint_L (\mathbf{v}\times \mathbf{B})\cdot \d \mathbf{l}
\]

我们考虑这样的一个例题:

\begin{example}
长度为$L$的一根铜棒, 绕其中一端在垂直于棒的均匀磁场中以角速度$\omega$旋转, 求这根铜棒两端的电势差.

铜棒旋转时切割磁感线, 故棒两端之间有感应电动势. 每个小段$\d l$产生的电动势方向相同, 于是我们可以计算:
\[
V = \mathscr C = \int_0^L B\omega l\,\d l =\frac{1}{2}B\omega L^2
\]

\end{example}
也许我们会有这样的问题, 注意到洛伦兹力对电荷永不做功, 这里我们又说动生电动势是由洛伦兹力做功引起的, 这岂非矛盾? 这种佯谬产生的原因是我们只计算了洛伦兹力的一部分. 在运动导体中的电子不仅有导体本身的速度, 还有相对导体的定向运动速度$u$, 正是这种运动构成了感应电流. 这个分量的方向恰好沿着$-\mb{v}$, 阻碍导体运动, 做负功阻碍导体的运动. 因此洛伦兹力只是将外力克服洛伦兹力的一个分量所做的功通过另一个分量转化为感应电流的能量.

导体在磁场中运动产生动生电动势, 其非静电力是洛伦兹力; 在磁场变化产生感生电动势的情形里非静电力是什么呢? 实验表明, 感生电动势完全与导体的种类和性质无关. 这说明感生电动势是由变化的磁场本身引起的.

我们考虑一个固定的回路$L$, $S$为以$L$为边界的曲面. $Maxwell$提出变化的磁场会产生涡旋电场$\mathbf{E}_i$. 产生感应电动势的非静电力就是涡旋电场, 这也提示我们静电场与涡旋电场并非同一种电场, 涡旋电场的环路积分不为$0$且恰好为感生电动势:
\[
\mathscr C = \oint_L \mathbf{E}_i \cdot  \d \mathbf{l} = -\frac{\d}{\d t} \iint_S \mathbf{B}\cdot \d \mathbf{S} = -\iint_S \frac{\partial \mathbf{B}}{\partial t}\cdot \d \mathbf{S}
\]
涡旋电场的存在是$Maxwell$电磁理论的基本假设之一.

现在我们来看一个感生电动势的应用---电子感应加速器. 它的基本结构如图所示.

\begin{figure}[h]
  \centering
  % 第一张图片
  \begin{minipage}[t]{0.3\textwidth}
    \centering
    \includegraphics[width=\textwidth]{4-2.png}
    \caption{电子感应加速器}
  \end{minipage}
  \hfill
  % 第二张图片
  \begin{minipage}[t]{0.3\textwidth}
    \centering
    \includegraphics[width=\textwidth]{4-3.png}
    \caption{感应加速器中磁场变化位于不同相位时涡旋电场的方向}
  \end{minipage}
\end{figure}

在一个电流变化的周期内, 只有四分之一的区间能够用于加速电子. 这是因为, 首先为了使电子得到加速, 涡旋电场应该沿着顺时针方向, 这要求第一或者第四个四分之一周期. 其次为了使得电子不断加速, 必须维持洛伦兹力指向圆心. 可以看出, 只有第一或者第二个四分之一周期. 综上所述, 只有第一个四分之一周期电子才能不断被加速.

为了实现这种加速, 我们需要同时维持电子的圆周运动和加速. 电子做圆形轨道时要求有向心力公式:
\[
evB = \frac{mv^2}{R} \Rightarrow mv = eBR \Rightarrow \d(mv) = eR\d B
\]

上式表明, 只要电子动量随着磁感应强度成比例增加, 就可以维持电子在轨道上运动. 电子的加速过程可以描述如下:
\[
E_{orb} \cdot 2\pi R = -\frac{\d \Phi}{\d t} \Rightarrow
E_{orb} = -\frac{1}{2\pi R}\frac{\d \Phi}{\d t}
\]
根据牛顿第二定律:
\[
\frac{\d (mv)}{\d t}=-eE_{orb}=\frac{e}{2\pi R}\frac{\d \Phi}{\d t} \Rightarrow \d(mv)  =\frac{e \d \Phi}{2\pi R}
\]
取$\overline{B}$为轨道内的平均磁感应强度, 则有$\d \Phi = \pi R^2\d \overline{B}$. 于是我们得到:
\[
\d B = \frac{1}{2}\d \overline{B} \Rightarrow B = \frac{1}{2}\overline{B}
\]
也就是说, 当轨道处磁感应强度为轨道内平均磁感应强度的一半时, 电子能够在稳定的轨道上被加速.

最后我们简单介绍趋肤效应. 当导线通入稳恒电流时, 横截面上的电流密度基本上是均匀分布的. 然而当导线通入交变电流时, 将会激发出磁场. 这一变化的磁场又会激发出涡旋电场, 从而产生涡电流. 在导线表面, 涡电流的方向与传导电流相同, 在靠近轴线处反之. 导线中的电流趋向于分布在导线的表面, 这一效应被称为趋肤效应.

\section{自感与互感}
\insertfig{4-4.png}{两线圈之间的互感}{0.3}
当线圈$1$中的电流变化时所激发的变化磁场, 会在它邻近的另一线圈$2$中产生感应电动势; 同样线圈$2$中的电流变化时也会在线圈$1$中产生感应电动势. 这种现象称为互感现象, 产生的感应电动势称为互感电动势. 显然一个线圈中的互感电动势不仅与另一线圈中的电流变化率有关, 而且也与两个线圈的结构以及它们之间的相对位置有关.

设线圈$1$所激发的磁场通过线圈$2$的磁链为$\Psi_{12}$按照毕奥-萨伐尔定律, $\Psi_{12}$与线圈$1$中的电流$I_1$成正比:
\[
\Psi_{12} = M_{12} I_1
\]
对称地, 我们有:
\[
\Psi_{21} = M_{21} I_2
\]

其中的$M_{12},M_{21}$被称为互感系数, 简称互感. 可以证明$M_{12}=M_{21}=M$. 当然这里我们讨论的过程中不涉及到铁磁质, 否则$M_{12},M_{21}$都不再是常数.

当电流$I_1$发生变化时, 将在线圈$2$中激发的感应电动势为:
\[
\mathscr C_2 = -\frac{\d \Psi_{12}}{\d t} = -M\frac{\d I_1}{\d t},\qquad \mathscr C_1 = -\frac{\d \Psi_{21}}{\d t} = -M\frac{\d I_2}{\d t}
\]
上面我们已经假定互感系数保持不变.

同样地, 当一个线圈中通入电流时, 线圈会产生磁场. 这一磁场对线圈自身是有磁通量的. 因此当线圈内部的电流发生变化时, 变化的磁场自然也会导致穿过线圈的磁通发生变化. 于是线圈中将产生感应电动势, 这一现象被称为自感现象. 一个线圈便可以称为一个自感.

我们知道线圈中电流所激发的磁感应强度与电流成正比, 因此通过线圈的磁链也正比于线圈中的电流, 即:
\[
\Psi = LI \Rightarrow \mathscr C = -L\frac{\d I}{\d t}
\]
$L$被称为自感系数. 同样地, 不存在铁磁质的情形下, 自感系数则为常数.

我们来考虑一个截面积为$S$, 长为$l$的$N$匝长直密绕螺线管的自感系数. 设该螺线管上的电流大小为$I$. 则其产生的磁感应强度大小为:
\[
B = \mu_0\frac{NI}{l}
\]
通过螺线管横截面的磁通量为$\Phi= BS=\mu_0\frac{NIS}{l}$. 于是自感系数为:
\[
L =\frac{N\Phi}{I}=\mu_0\frac{N^2S}{l} = \mu_0 n^2 V
\]

两个线圈之间的互感系数与各自的自感系数有一定的联系. 当两个线圈中每一个线圈所产生的磁通量对于每一匝来说都相等, 并且全部穿过另一个线圈的每一匝. 这种情况就被称为无漏磁. 一般而言, 通过自身回路的磁通量要大些.
\[
\Phi_{21} = k_1 \Phi_1, \;\Phi_{12} = k_2 \Phi_2, k_1\le 1, k_2\le 1
\]

考虑到先前的定义, 我们可以把磁链写成:
\[
\Psi_{21} = N_2 \Phi_{21} = N_2 k_1 \Phi_1,\; \Psi_{12} = N_1 \Phi_{12} = N_1 k_2 \Phi_2
\]

因此互感系数满足:
\[
M_{21}=\frac{N_2k_1\Phi_1}{I_1},\;M_{12}=\frac{N_1k_2\Phi_2}{I_2}\Rightarrow M^2 =\frac{N_1N_2k_1k_2\Phi_1\Phi_2}{I_1I_2}
\]

考虑自感关系, 有:
\[
\Psi_1 = N_1 \Phi_1 = L_1 I_1,\;\Psi_2 = N_2 \Phi_2 = L_2 I_2
\]

把上式代入, 得到:
\[
M = \sqrt{k_1k_2}\sqrt{L_1L_2}
\]

$k=\sqrt{k_1k_2}$被称为线圈$1$和线圈$2$之间的耦合系数.

若不存在漏磁, 则$k=1$:
\[
M=\sqrt{L_1L_2}
\]

\section{\texorpdfstring{暂态过程\, 磁场能量}{暂态过程 磁场能量}}
暂态过程指的是是电路状态从一种稳态转变为另一种稳态时所经历的过渡过程. 这种暂态过程通常在电路发生突变, 比如开关闭合, 电源接入, 初始条件发生变化时. 电路状态改变时, 电容, 电感等元件需要时间来调整其储存的能量. 这就产生了暂态过程.

我们来看一个电路, 将所有的电容, 电阻, 电感分别等效为$C,R,L$三个基本元件, 设初态$t=0$时电路中$I=0$, 电容极板上电荷$q=0$. 于是我们可以写出回路中的电势降方程:
\[
L\frac{\d I}{\d t}+\frac{q}{C}+IR = \varepsilon
\]
注意到:$I=\frac{\d q}{\d t}$, 可以整理我们的方程为:
\[
L\frac{\d^2 q}{\d t^2}+R\frac{\d q}{\d t}+\frac{q}{C}=\varepsilon
\]

我们讨论一下几种特殊情况的电路:

(1)$L=0$. 这是一个$RC$电路. 解上述方程, 可以得到:
\[
q =\v C(1-e^{-\frac{t}{RC}}),\; I=\frac{\v}{R} e^{-\frac{t}{RC}}
\]
其中$\tau=RC$被称为特征时间, 可以用来描述电路达到稳态的时间长短. 我们可以看到, 在充电的过程中, 电容上的电量逐步从$0$开始增加, 最终达到饱和. 而电流则以指数衰减为$0$.

放电时, 取$\v=0$电量和电流大小都以指数形式衰减:
\[
q = q_0e^{-\frac{t}{RC}},\; I = -\frac{q_0}{RC}e^{-\frac{t}{RC}}
\]
其中$q_0$是初始状态下电容上的电量.

(2)移去电容器. 这是一个$RL$电路, 方程应直接写为$L\d I/\d t+RI=\v$. 解上述方程, 可以得到:
\[
I =\frac{\v}{R}(1-e^{-\frac{R}{L}t})
\]
其中$\tau=\frac{L}{R}$为特征时间.

这个解表明, 在初始状态下, 电流为$0$. 后电流逐步以指数增长形式达到稳态. 放电过程与之恰好相反:
\[
I = I_0 e^{-\frac{R}{L}t}
\]

在以上的讨论中, 我们发现电容上的电压和电荷不能瞬间变化, 电感中的电流不能瞬间变化. 否则会产生无穷大电动势.

(3)$R$=0.这是一个$LC$电路. 解上述方程, 我们得到:
\[
q=q_0\cos(\omega_0 t+\phi),\quad \omega_0 = \frac{1}{\sqrt{LC}}
\]
也就是说, 在这种电路中电荷和电流都将以余弦形式震荡, 形成我们后面会讨论的产生电磁波的物质基础.

(4)一般情况. 这时候, 我们把整个方程化为标准的二阶常系数微分方程:
\[
\frac{\d^2 q}{\d t^2}+2\beta \frac{\d q}{\d t}+\omega_0^2 q = \frac{\v}{L},\quad \beta = \frac{R}{2L}, \quad \omega_0 = \frac{1}{\sqrt{LC}}
\]

这方程的通解与$RLC$的具体数值有关, 我们可以定义阻尼度:
\[
\lambda = \frac{R}{2}\sqrt{\frac{C}{L}}
\]

$\lambda<1$为欠阻尼, $\lambda=1$为临界阻尼, $\lambda>1$为过阻尼. 欠阻尼情况下, 电路的振荡频率为:
$\omega = \sqrt{\frac{1}{LC}-(\frac{R}{2L})^2}$. 上述方程的一个特解是$q=\v C$.

讨论了暂态过程, 我们现在来看一下磁场中的能量. 我们考虑在一个线圈中建立电流的过程. 当线圈与电源接通时, 由于自感的存在的, 电路中的电流$i$并不立刻变化到稳定值$I$. 在这一变化过程中, 电路中外电源的电动势不仅要供给电路中的焦耳热, 而且要反抗自感电动势的做功, 这部分反抗自感电动势所做的额外的功应当以能量的形式储存在线圈中.

$\d t$时间内, 电源的电动势反抗自感电动势做功为:
\[
\d A = -\mathscr C_{L} i \d t =-(-L\frac{\d i}{\d t})i\d t = Li\d i
\]
因而在整个建立电流的过程中, 电源的电动势反抗自感电动势所做的功为:
\[
A = \int_0^I Li\d i =\frac{1}{2}LI^2
\]
当切断电源时, 电流由稳定值$I$减少到$0$. 线圈中产生了阻碍电流减小的自感电动势, 其驱动方向与原电流方向相同. 线圈中原本存储的能量将会以同样的形式释放出来. 于是我们便说, 自感线圈将会存储磁能, 在一个自感系数为$L$的线圈中建立强度为$I$的电流, 线圈中存储的能量是$\frac{1}{2}LI^2$. 这被称为自感磁能.

下面简单讨论互感磁能: 我们考察在相邻的线圈$1$和$2$中建立电流的过程中抵抗互感电动势所做的功:

\[
A = -\int_0^{\infty}\mathscr C_{21}i_1\d t-\int_0^{\infty}\mathscr C_{12}i_2\d t =\int_0^{\infty} M(i_1\frac{\d i_2}{\d t}+i_2\frac{\d i_1}{\d t})\,\d t=M\int_0^{\infty}\frac{\d(i_1i_2)}{\d t} \d t = MI_1I_2
\]
也就是说线圈中存储的互感能量为$MI_1I_2$.

我们从螺线管入手, 导出磁场能量的一般形式: 考虑螺线管的磁导率为$\mu$, 长度为$l$, 截面积为$S$, 线圈匝数为$N$, 电流强度为$I$, 则环内磁场为:
\[
B  =\mu nI
\]
螺线管的自感系数在之前已经推导过:
\[
L = \frac{\mu N^2S}{l} = \mu n^2 V
\]
于是螺线管储存的磁能为:
\[
W_m = \frac{1}{2}LI^2 = \frac{1}{2}\mu n^2 I^2 V = \frac{1}{2}VBH
\]
于是螺线管内的磁能密度为:
\[
w_m = \frac{1}{2}\mathbf{B}\cdot \mathbf{H}
\]
上式适用于线性、无磁滞且参数不随时间变化的介质; 对非线性磁介质, 应从$w_m=\int_0^{\mathbf B}\mathbf H\cdot\d\mathbf B$求能量密度.

